\chapter{Constraint Manifolds}

\section{What This Part Owes the Reader}

Six chapters of Part III were written against a promise this book has not yet kept. C appeared in Chapter 14 as the constraint manifold, the set of admissible field states, and has been invoked in every chapter since: \(\Pi_C\) projected proposals onto it, argmin measured distance to it, \(M^*\) solved an equation stated in terms of it. At no point was C actually defined as a mathematical object, only gestured at through the rules that were supposed to carve it out. This was a deliberate deferral, consistent with the four-pass structure this book announced in its roadmap: Parts II and III developed C's role architecturally, establishing what it needed to do, before Part IV develops what it actually is. This chapter keeps that promise, and the chapters following it keep the promises Part IV itself will go on to make.

\section{Assembling C from Its Members}

The honest way to define C is not to write down a single equation and declare it the constraint manifold, but to collect what earlier chapters already established and show that their intersection is what C has been all along.

Chapter 8 contributed the entropy constraint: the set of states for which \(\Delta S_t(x)\) ≥ 0 at every point outside the region actually observed during the update in question. Call this \(C_S\). Chapter 9 contributed appearance consistency: the set of states in which a settled \(z_t(x)\) is not overwritten by a fresh candidate without cause, formalized through the invention-to-memory update itself. Call this \(C_z\). Chapter 11 contributed residue consistency: the set of states whose appearance and entropy values do not contradict what accumulated residue at that point has already ruled out. Call this \(C_R\). Chapter 10 promised, and Chapter 21 will formalize, a further member governing the realizability of affordance-driven proposals as genuine field transformations; call this \(C_A\), understood for now only as a placeholder this chapter cannot yet complete.

C is the intersection of these:

\[ C = C_S \cap C_z \cap C_R \cap C_A \cap \cdots, \]

and the ellipsis is honest rather than decorative. Nothing in this book's argument so far guarantees that the list is complete, and Part IV's later chapters, on conservation laws, category theory, and information geometry, will each have occasion to add further members before the list can be called closed. What matters for this chapter is only that C has never been an arbitrary constraint invented for Part IV's convenience. It is the accumulated weight of every discipline Parts II and III already argued for, now written as a single intersection rather than scattered across six chapters.

\section{The Space the Field Actually Lives In}

For "nearest point," "intersection," and "manifold" to mean anything precise rather than serving as evocative language borrowed from geometry, the space these terms operate in has to be specified. This book takes the field's configuration space to be a genuine differentiable manifold, written 𝓜, equipped with a metric g that gives the norm in Chapter 14's projection equation, ‖Y − (\(M_t\) + \(\Delta\psi\))‖², actual mathematical content rather than standing in as a placeholder for "some notion of closeness." A point of 𝓜 is a complete assignment of Φ, v, S, R, and z across the field's domain X; the metric g measures how far apart two such assignments are, in whatever combination of coherence-distance, flow-distance, entropy-distance, residue-distance, and appearance-distance the full theory eventually specifies component by component.

C, then, is not a separate kind of object from 𝓜. It is a subset of 𝓜, carved out by the intersection of conditions in section 18.2, and in the well-behaved cases this book will generally assume, a submanifold: a lower-dimensional surface sitting inside the larger space of all conceivable, not-necessarily-admissible field configurations. Projection, \(\Pi_C\), is now exactly what differential geometry already means by the term: given a point of 𝓜 not lying on the submanifold C, find the nearest point of C under the metric g. Nothing about Chapter 14's equation was metaphorical. It was waiting for this section to specify the space in which it is literally true.

\section{Hard Constraints and Soft Dynamics}

It is worth separating, within C's definition, two kinds of condition that Part III's supporting material already distinguished informally. Hard constraints fix C's shape outright: mass consistency, identity persistence, causal locality, and topological admissibility are not negotiable by any proposal, however well motivated, and no \(\Delta\psi\), however strongly selected by the relevance mechanism of Chapter 13, can move the field outside the region these conditions carve out. Soft dynamics, by contrast, are what varies within whatever region the hard constraints permit: the specific deformation a courtyard's coherence undergoes, the specific texture an appearance settles into, the specific trajectory a proposed affordance realization follows. \(F_{\mathrm{phys}}\) and \(F_\psi\), introduced in Chapter 13, were always proposing within the space hard constraints define, never proposing changes to the constraints themselves. This distinction gives precise content to a claim Chapters 12 through 17 made repeatedly without formalizing: that physics, in this book's sense of the word, lives in C, while the write cycle's proposal machinery lives in the freedom C leaves open.

\section{When the Domain Itself Changes}

Chapter 8 stated its entropy constraint as unconditional, and deliberately declined to qualify it, on the grounds that qualifying it there would have weakened a chapter whose entire force depended on stating a rule plainly. That deferral is settled here. The entropy constraint, and the hard constraints generally, are unconditional on a fixed field domain X. They say nothing, on their own, about what happens when X itself changes: when a topology repair merges two previously distinct regions, or an identity restructuring splits one region into two, altering the very space over which Φ, v, S, R, and z are defined.

The resolution is transport rather than exception. Given a structure-preserving map φ relating the old domain to the new one, whatever hard constraints held on the old domain are carried across by φ and continue to hold on the new one: a condition stated as \(C_S\) on X becomes \(\varphi_*C_S\), its pushforward, on the new domain X′, and the entropy constraint on X′ is exactly this transported condition rather than a fresh rule reasoned out from nothing. Nothing about the constraint has been relaxed, and nothing about it needed to be. What has changed is only the space it is being asked to hold over, and transport is precisely the operation that lets a rule stated once continue to mean the same thing after the space beneath it has moved. This book will not develop the full apparatus governing which maps φ are eligible to carry constraints across a topology change; that question belongs with the sheaf-theoretic material Chapter 21 develops, since a topology change of this kind is exactly the situation in which local field-patches need to be reconciled with a new global cover.

\section{What This Chapter Has Not Established}

It is worth being explicit about the gap this chapter leaves open, because Chapter 16 asked a question this chapter's machinery still cannot answer. \(M^*\) = \(\Pi_C\)(\(M^*\) + \(\Delta\psi\)) presumes that a nearest point in C exists for whatever is being projected, and existence of a nearest point is not automatic once C is merely known to be a submanifold. A submanifold that is not closed, or a metric that is not complete, can leave some proposals with no nearest admissible point at all, or with several equally near, and Chapter 14's argmin silently assumed that this pathology would not arise. Nothing established in this chapter rules it out. Resolving it requires more than a definition of C; it requires understanding the dynamics that move a proposal toward C in the first place, and whether those dynamics are governed by a well-behaved energy functional whose minimization is guaranteed to terminate somewhere. That is a variational question, not a purely geometric one, and it is where Chapter 19 begins.

\section{Looking Ahead}

This chapter has given C a real definition: the intersection of every discipline Parts II and III argued for, sitting as a submanifold inside the field's full configuration space, distinguished into hard constraints that fix its shape and soft dynamics that vary within it, with a transport rule now on record for what happens when the underlying domain itself changes. What it has not given is any guarantee that projection onto C behaves well, that a nearest point exists, or that the write cycle's repeated application of \(\Pi_C\) actually settles into the fixed points Chapter 16 described rather than oscillating or failing to converge at all. Chapter 19 supplies the missing piece: an energy functional, in the proper variational sense, whose minimization is what argmin was always secretly asking for.
